
%(BEGIN_QUESTION)
% Copyright 2006, Tony R. Kuphaldt, released under the Creative Commons Attribution License (v 1.0)
% This means you may do almost anything with this work of mine, so long as you give me proper credit

Hvorfor er \textit{temperatur} en viktig faktor i trykktransmittere med remote seals (ofte kalt ``chemical seals'')?  Hvilke temperaturforhold kan gi målefeil i et differansetrykkinstrument med remote seals?

\vskip 10pt

Beskriv også hva som skjer med et instrument med remote seals dersom det oppstår lekkasje i ett av kapillarrørene.

\underbar{file i00217}
%(END_QUESTION)





%(BEGIN_ANSWER)

Økt temperatur får fill-væsken i kapillarrøret til å utvide seg, og dette skaper et trykk på transmitterens sensorelement uavhengig av prosessens reelle trykk.

Hvis instrumentet er differensielt og begge seal/kapillarsider har samme temperatur, vil temperaturinduserte trykk i stor grad kansellere.  Hvis temperaturen derimot er ulik mellom sidene, oppstår et \textit{differansetrykk} fra temperatur alene, som transmitteren kan tolke som prosessdifferansetrykk.

\vskip 10pt

Remote seal/kapillarsystemer må være helt gassfrie (kun væske) for å fungere riktig.  Selv en liten luft/gassboble i fill-væsken gjør systemet komprimerbart.  Da vil seal-membranen flytte et større væskevolum enn sensorelementet, og instrumentet ``ser'' ikke hele prosessendringen.  Målenøyaktigheten forringes.

Kort oppsummert: lekkasjer i remote seal-systemer er svært skadelige for nøyaktigheten.  Hvis lekkasje oppstår, må hele instrumentet normalt byttes.  Selv ved reparert lekkasjepunkt må systemet vakuumfylles på nytt for å trekke ut oppløste/innestengte gasser, noe som sjelden kan utføres i vanlig instrumentverksted.

%(END_ANSWER)





%(BEGIN_NOTES)


%INDEX% Measurement, pressure: chemical seals
%INDEX% Measurement, pressure: remote seals

%(END_NOTES)

